\section{Introduction}
\label{sec:introduction}

What if the key to safer AI isn't teaching models what \emph{not} to do, but teaching them to follow any constraint---including implicit ones they've never seen?

Reinforcement Learning from Human Feedback (RLHF) has become the dominant paradigm for aligning language models with human preferences \citep{ouyang2022instructgpt}. Current approaches optimize a single objective: helpfulness as judged by human annotators or proxy reward models. This \emph{unidirectional} alignment (AI$\rightarrow$Human) has produced increasingly capable assistants, yet a fundamental gap remains: models follow instructions inconsistently, particularly explicit constraints on format, length, and structure \citep{zhou2023ifeval}.

Sun et al.\ \citep{sun2024bidirectional} propose a \emph{bidirectional} alignment framework distinguishing two directions: AI$\rightarrow$Human (helpfulness, harmlessness) and Human$\rightarrow$AI (controllability, steerability). While the theoretical framework is compelling, no existing work implements bidirectional training signals. IFEval \citep{zhou2023ifeval} measures instruction-following for evaluation; AlpacaEval \citep{dubois2024alpacaeval} measures helpfulness---but these metrics remain siloed, each used only for post-hoc evaluation rather than as integrated training signals.

We address this gap with a simple intervention: augmenting the standard helpfulness reward with an IFEval-derived controllability signal. Our combined reward takes the form:
\begin{equation}
R = \alpha \cdot R_{\text{helpfulness}} + \beta \cdot R_{\text{IFEval}}
\end{equation}
where $\alpha, \beta$ control the trade-off between objectives. We convert discrete IFEval constraint checks into continuous, differentiable rewards via sigmoid soft thresholds, enabling gradient-based optimization.

Our experiments validate three predictions:

\begin{enumerate}
    \item \textbf{P1 (Controllability):} Bidirectional models achieve higher held-out IFEval accuracy than helpfulness-only baselines (+3.4pp strict accuracy).
    \item \textbf{P2 (Helpfulness Maintenance):} At $\alpha \geq 0.8$, bidirectional models retain $\geq$95\% of baseline AlpacaEval performance (96.4\% observed).
    \item \textbf{P3 (Safety Transfer):} Explicit constraint training transfers to implicit safety constraints, with 2--4pp gains on TruthfulQA and BBQ and strong positive correlation ($r=0.85$--$0.94$) between IFEval improvement and safety improvement.
\end{enumerate}

The transfer finding (P3) is our most surprising result. Models trained to follow explicit format constraints---``respond in exactly 3 paragraphs,'' ``include the word `conclusion'\,''---show improved performance on safety benchmarks measuring truthfulness and bias. This suggests that controllability training builds general constraint-following capacity that extends beyond the explicit constraints seen during training.

\paragraph{Contributions.} We provide:
\begin{itemize}
    \item The first empirical test of bidirectional alignment training signals, validating the theoretical framework of Sun et al.\ \citep{sun2024bidirectional}
    \item A methodology for repurposing rule-based evaluation metrics (IFEval) as differentiable RLHF training rewards
    \item A quantified estimate of explicit$\rightarrow$implicit transfer: $\sim$15\% of IFEval improvement transfers to safety metrics
    \item Pareto characterization of the helpfulness-controllability trade-off, identifying viable operating points for different deployment priorities
\end{itemize}
